\documentclass[12pt,a4paper]{article}
\usepackage{amsmath,amssymb,amsfonts}
\usepackage{graphicx}
\usepackage{hyperref}
\usepackage{physics}
\usepackage{siunitx}
\usepackage{booktabs}
\usepackage{bm}

\title{Simulation of the Berezinskii-Kosterlitz-Thouless Transition in the 2D XY Model}
\author{CSP Project Report}
\date{\today}

\begin{document}
\maketitle

\tableofcontents
\vspace{1cm}

\section{Introduction}

The two-dimensional XY model exhibits a peculiar phase transition known as the Berezinskii-Kosterlitz-Thouless (BKT) transition. Unlike conventional phase transitions characterized by spontaneous symmetry breaking and a non-zero order parameter, the BKT transition represents a topological phase transition where the system transitions from a phase dominated by bound vortex-antivortex pairs to a phase where these pairs dissociate. This report describes a Monte Carlo simulation implementation that studies this transition using advanced sampling techniques.

The Monte Carlo code simulates the 2D XY model on square lattices of different sizes to observe finite-size scaling behavior near the critical temperature. The simulation employs the Wolff cluster algorithm for efficient sampling, particularly effective near the critical point, combined with parallel tempering (replica exchange) to improve ergodicity.

\section{Theoretical Background}

\subsection{The XY Model Hamiltonian}

The 2D XY model consists of planar rotors (spins constrained to move in a plane) arranged on a square lattice. Each spin is represented by a unit vector $\mathbf{S}_i = (\cos\theta_i, \sin\theta_i)$ characterized by an angle $\theta_i$ that can take any value in $[0, 2\pi)$. The Hamiltonian of the system is:

\begin{equation}
H = -J \sum_{\langle i,j \rangle} \mathbf{S}_i \cdot \mathbf{S}_j = -J \sum_{\langle i,j \rangle} \cos(\theta_i - \theta_j)
\end{equation}

where $J > 0$ is the ferromagnetic coupling constant, and the sum runs over nearest-neighbor pairs $\langle i,j \rangle$. 

\subsection{The Berezinskii-Kosterlitz-Thouless Transition}

The BKT transition occurs at a critical temperature $T_{KT} \approx 0.89 J/k_B$. Below $T_{KT}$, the system exhibits quasi-long-range order characterized by algebraically decaying correlation functions. Above $T_{KT}$, correlations decay exponentially.

The critical behavior is associated with topological defects called vortices. A vortex is a point around which the spins rotate by $2\pi$. The BKT transition can be understood as a transition from bound vortex-antivortex pairs at low temperatures to free vortices at high temperatures.

A key prediction of BKT theory is the universal jump in the spin stiffness (helicity modulus) at the transition:

\begin{equation}
\rho_s(T_{KT}) = \frac{2T_{KT}}{\pi}
\end{equation}

This relation allows us to determine $T_{KT}$ by finding where $\rho_s$ intersects the line $2T/\pi$.

\section{Monte Carlo Methods}

\subsection{Metropolis Algorithm}

The classical Metropolis algorithm updates spins one at a time:
\begin{enumerate}
    \item Select a random spin $i$ with angle $\theta_i$
    \item Propose a new angle $\theta_i' \in [0, 2\pi)$
    \item Calculate the energy change $\Delta E = E_{\text{new}} - E_{\text{old}}$
    \item Accept the new state with probability $P = \min(1, \exp(-\beta \Delta E))$
\end{enumerate}

However, near the critical point, the Metropolis algorithm suffers from critical slowing down, making it inefficient.

\subsection{Wolff Cluster Algorithm}

The Wolff algorithm addresses critical slowing down by updating entire clusters of spins:
\begin{enumerate}
    \item Choose a random reflection axis perpendicular to a random unit vector $\mathbf{r}$
    \item Pick a random site as the initial cluster
    \item Grow the cluster by adding neighboring spins with probability $P = 1 - \exp(-2\beta J \mathbf{S}_i \cdot \mathbf{S}_j)$ if they haven't been visited
    \item Reflect all spins in the cluster with respect to the chosen axis
\end{enumerate}

In our implementation, the reflection uses a random angle $\theta_{\text{rand}}$ and the transformation $\theta_i \rightarrow 2\theta_{\text{rand}} - \theta_i$ for all spins in the cluster.

\subsection{Parallel Tempering}

To further improve sampling efficiency, the simulation implements parallel tempering (replica exchange):
\begin{enumerate}
    \item Run simulations at multiple temperatures simultaneously
    \item Periodically propose exchanges between configurations at adjacent temperatures
    \item Accept exchanges with probability $P = \min(1, \exp(\Delta\beta \Delta E))$
\end{enumerate}

This allows configurations to move between high and low temperatures, preventing them from getting trapped in local energy minima.

\section{Implementation Details}

\subsection{Code Structure}

The simulation code is organized into several components:
\begin{itemize}
    \item Core Monte Carlo functions implementing the Wolff and Metropolis algorithms
    \item Data analysis functions for calculating thermodynamic observables
    \item Statistical analysis using jackknife resampling to estimate errors
    \item Plotting functions to visualize results and perform finite-size scaling
\end{itemize}

The simulation follows these steps:
\begin{enumerate}
    \item Pre-thermalization to determine optimal number of Wolff updates
    \item Thermalization to reach equilibrium
    \item Measurement phase to collect data
    \item Analysis of results to calculate observables and errors
\end{enumerate}

\subsection{Optimization Techniques}

Several optimizations are employed to improve performance:
\begin{itemize}
    \item Just-in-time compilation using Numba for compute-intensive functions
    \item Dynamic adjustment of the number of Wolff updates based on average cluster size
    \item Pre-computation of neighbor lists to avoid redundant calculations
    \item Parallel tempering to enhance ergodicity and convergence
\end{itemize}

\section{Physical Observables}

The simulation measures several key observables:

\subsection{Energy and Specific Heat}

The energy per spin is:
\begin{equation}
e = \frac{1}{N}\langle H \rangle = -\frac{J}{N} \sum_{\langle i,j \rangle} \langle\cos(\theta_i - \theta_j)\rangle
\end{equation}

The specific heat per spin is calculated from energy fluctuations:
\begin{equation}
c_v = \frac{\beta^2}{N}(\langle H^2 \rangle - \langle H \rangle^2) = \frac{1}{NT^2}(\langle E^2 \rangle - \langle E \rangle^2)
\end{equation}

\subsection{Magnetization and Susceptibility}

The magnetization per spin is:
\begin{equation}
\mathbf{m} = \frac{1}{N}\left\langle \sum_i \mathbf{S}_i \right\rangle = \frac{1}{N}\left\langle \sum_i (\cos\theta_i, \sin\theta_i) \right\rangle
\end{equation}

The magnetic susceptibility per spin is:
\begin{equation}
\chi = \frac{\beta}{N}(\langle M^2 \rangle - \langle M \rangle^2)
\end{equation}

\subsection{Spin Stiffness (Helicity Modulus)}

The spin stiffness measures the system's resistance to a twist in boundary conditions:
\begin{equation}
\rho_s = \left\langle \frac{1}{N} \sum_{\langle i,j \rangle} \cos(\theta_i - \theta_j) \right\rangle - \frac{\beta}{N} \left\langle \left(\sum_{\langle i,j \rangle} \sin(\theta_i - \theta_j)\right)^2 \right\rangle
\end{equation}

This is a crucial quantity for identifying the BKT transition.

\subsection{Vortex Density}

The vortex density is calculated by counting the net winding number of the spin orientations around each plaquette:
\begin{equation}
\omega_v = \frac{1}{N} \sum_{\text{plaquettes}} \left| \frac{1}{2\pi} \sum_{\text{edges}} \Delta\theta_{ij} \right|
\end{equation}
where $\Delta\theta_{ij}$ is the principal value of the angle difference between neighboring spins, constrained to $[-\pi, \pi]$.

\section{Error Analysis}

To estimate statistical errors, the simulation employs the jackknife resampling method:
\begin{enumerate}
    \item Divide data into $N_B$ blocks of length $k$
    \item Calculate block averages
    \item For each block $i$, compute the jackknife estimate by excluding block $i$
    \item Calculate the mean and variance of these jackknife estimates to determine statistical errors
\end{enumerate}

This approach provides more reliable error estimates than simple variance calculations, especially for derived quantities.

\section{Finite-Size Scaling Analysis}

To extract the thermodynamic limit of $T_{KT}$, we perform finite-size scaling analysis. The BKT theory predicts that $T_{KT}(L)$ scales with system size as:
\begin{equation}
T_{KT}(L) = T_{KT}(\infty) + \frac{a}{(\ln L)^2}
\end{equation}

By plotting $T_{KT}(L)$ against $1/(\ln L)^2$ for different system sizes and performing a linear fit, we can extrapolate to find $T_{KT}(\infty)$.

\section{Results and Interpretation}

The simulation results are presented in several plots:

\subsection{Specific Heat}

The specific heat exhibits a broad peak that sharpens slightly with increasing system size. Unlike conventional phase transitions, the specific heat peak does not diverge at the BKT transition.

\subsection{Spin Stiffness}

The spin stiffness shows a characteristic behavior near the transition temperature. The intersection of $\rho_s(T)$ with the universal jump line $2T/\pi$ provides an estimate of $T_{KT}$ for each system size.

\subsection{Vortex Density}

The vortex density increases rapidly near the transition temperature as vortex-antivortex pairs unbind. This provides a direct visualization of the topological nature of the transition.

\subsection{Susceptibility}

The magnetic susceptibility shows a diverging behavior with system size, reflecting the establishment of quasi-long-range order below $T_{KT}$.

\section{Conclusion}

The Monte Carlo simulation successfully captures the essential features of the BKT transition in the 2D XY model. The finite-size scaling analysis provides an estimate of $T_{KT}$ consistent with theoretical predictions and previous numerical studies. The implementation demonstrates the power of advanced Monte Carlo techniques such as cluster algorithms and parallel tempering in studying complex phase transitions.

The transition is most clearly identified through the behavior of the spin stiffness and its intersection with the universal jump line. The results confirm the topological nature of the BKT transition, characterized by vortex-antivortex pair unbinding rather than conventional symmetry breaking.

\end{document}
